\section{Introduction}

A wide class of latency-critical industrial problems reduces to the same computational
primitive: given a large sparse graph $G = (V, E)$ stored in memory and a source
vertex $s$, find all vertices reachable within $k$ hops and the shortest-hop distance
to each. This primitive underlies:

\begin{itemize}[leftmargin=*,itemsep=2pt]
\item \textbf{Financial fraud and risk detection}: identify transaction ring-members,
  card-sharing clusters, and account compromise chains in real time~\cite{akoglu2015graph}.
\item \textbf{Graph databases}: answer neighborhood queries, path existence tests,
  and reachability checks over entity-relationship graphs~\cite{besta2019demystifying}.
\item \textbf{Genomics}: traverse de Bruijn assembly graphs for contig construction
  in DNA sequencing pipelines~\cite{pevzner2001eulerian}.
\item \textbf{Recommendation}: multi-hop collaborative filtering over user-item bipartite
  graphs~\cite{ying2018graph}.
\end{itemize}

Each vertical demands different performance contracts, but all share the same memory
access signature: highly irregular, pointer-chasing reads that defeat CPU caches and
GPU batch pipelines alike. A k=3 BFS over a 1-million-vertex power-law graph requires
visiting $O(10^4$--$10^5)$ distinct vertices, each residing at a random location in a
tens-of-gigabyte memory space. The row-buffer hit rate $\lambda$ for such access
patterns is below 0.10, compared to 0.88 for regular stencil codes~\cite{fritz2026rmwwall}.

\subsection{The Hardware Gap}

Table~\ref{tab:competitive} summarizes the latency and power landscape for k=3 BFS
on a 1-million-vertex graph.

\begin{table}[h]
\centering
\caption{Competitive landscape: k=3 BFS, 1M vertices.}
\label{tab:competitive}
\small
\begin{tabular}{@{}lrr@{}}
\toprule
Platform           & Latency        & Power \\
\midrule
CPU (server, LLC)  & 50--500~$\mu$s & $\sim$200~W \\
GPU (A100, HBM2e)  & 10--50~$\mu$s  & $\sim$400~W \\
FPGA (Alveo U280)  & 2--10~$\mu$s   & $\sim$75~W  \\
\textbf{Target ASIC} & \textbf{0.5--2~$\mu$s} & \textbf{$\sim$5~W} \\
\bottomrule
\end{tabular}
\vspace{3pt}
\begin{minipage}{\linewidth}
\footnotesize\textit{CPU latency estimated from Graph500 published runs on server
hardware~\cite{graph500}. GPU from Merrill et al.\ scalable BFS on A100~\cite{merrill2012scalable},
adjusted for kernel launch overhead. FPGA from Betkaoui et al.~\cite{betkaoui2012reconfigurable}
and Oguntebi and Olukotun~\cite{oguntebi2016graphops}, scaled to HBM2e bandwidth.
ASIC is a pre-silicon design target.}
\end{minipage}
\end{table}

CPU performance is limited by cache miss chains: on a random graph, a 3-hop traversal
incurs $O(D_\text{avg}^k)$ LLC misses, each paying DRAM round-trip latency.
GPU performance is limited by kernel dispatch overhead ($\sim$5~$\mu$s minimum) and the
requirement to batch queries for effective SM utilization.
Neither platform provides deterministic tail latency, which is required for real-time
fraud scoring and HFT risk systems.

A purpose-built ASIC eliminates both penalties. Fixed-function hardware can stream CSR
adjacency data from HBM with no dispatch overhead, pipelining read requests to maximize
bus utilization, and providing hard real-time guarantees on traversal latency.

\subsection{A Cybernetic Framing}

The BFS traversal loop is a cybernetic system in Wiener's sense~\cite{wiener1948cybernetics}:
a feedback control loop that uses information received from the environment (HBM adjacency
lists) to update system state (the visited bitmap) and direct future action (the frontier
queue). The $\eta \approx 40\%$ efficiency bound established in~\cite{fritz2026rmwwall}
quantifies the communication rate between system and environment; it is the rate at which
the system can reduce uncertainty about the reachable subgraph. Custom silicon improves
this not by widening the channel, but by reducing control loop overhead --- the cycles
spent on protocol rather than information.

This framing matters beyond the immediate architecture. Mapping problems onto the
available technology is one step in a longer journey. The deeper question is what
becomes possible when the control loop closes fast enough that the machine's behavior
becomes indistinguishable from purpose. BFS-HBM is an existence proof at the memory
interface level. The same principle scales: as loop closure time falls, the class of
problems that can be solved in real time expands. The ASIC is not the end of the
road --- it is a measurement of where the road is now.

\subsection{Contributions}

This paper presents BFS-HBM, a fully synthesizable open-source RTL implementation of
the BFS traversal primitive for HBM-attached graphs:

\begin{enumerate}[leftmargin=*,itemsep=2pt]
\item A 12-state AXI4 traversal FSM (\texttt{bfs\_ctrl}) that exploits CSR memory
  layout to fetch both row-pointer endpoints in a single HBM transaction.
\item A 2-cycle BRAM read-modify-write visited bitmap (\texttt{visited\_bitmap})
  that tracks visited state for up to $2^{20}$ vertices in 128~KB of on-chip SRAM.
\item A cycle-accurate verification of correct BFS ordering against a golden reference,
  with a modeled HBM latency of 20 cycles at 250~MHz.
\item A closed-form throughput model connecting the RTL microarchitecture directly
  to the $\eta \approx 40\%$ HBM efficiency result of the companion bandwidth analysis.
\item A clear path from this RTL baseline to FPGA validation on Xilinx Alveo U280
  and ASIC tapeout on TSMC N7 or GlobalFoundries 12LP.
\end{enumerate}

All RTL, testbenches, and synthesis scripts are available at
\url{https://github.com/quantumcelnav/bfs-hbm-engine} under the MIT license.
